\section{Interfaces with singular geometry and dynamics}\label{sec:interfaces}
The foundations established here do not require all experiments to be
contractions of one physical path law or one rate function. They give a
common language in which an actual map or causal simulator must be exhibited
before a comparison is asserted.

\subsection{What is established in this installment}
The exact results are path construction, future-test minimality with stated
realization conditions, charged simulation composition, posterior and
entropy identities, the finite/dominated score identity, normalized
response, product-observation fibres, checkpoint quantization, and
exponential composition. The quantitative results are the feedback- and
stopping-compatible finite-state emulation theorem and the uniform dynamic
memory law for the refreshing model. Their proofs occur in this article.
The LDP and operator discussions are conditional interfaces, with a standard
contraction principle and a finite-dimensional memory identity; no general
Sinai or kinetic limit theorem is asserted here.

The principal dynamic model verifies a compact shrinking domain, actual
probability mass, and stable updates directly. It gives one time-uniform
isotropic degeneration law. It does not classify all anisotropic
collisions, all hidden Markov filters, all priors, or weak-refreshing limits.
The bounded comparison constants are not exact optimizer transition
constants.

\subsection{The role of the two geometric manuscripts}
The finite-horizon manuscript \cite{QiA1} provides finer anisotropic
profiles for its positive monomial experiments and a separate bounded
observation-algebra law. Those hypotheses should remain attached to their
conclusions. The present dynamic result is not a replacement for their
collision geometry, and their fixed-horizon constants are not assumed
uniform as the horizon grows.

The contact manuscript \cite{QiA2} provides a different geometric
realization of information loss. Theorem~\ref{thm:product-fibre} explains
its linear-functional interface. A complete extension to singular
statistical resolution still requires the actual noise metric, calibration,
parameter nuisance, and admissible real domain. In particular independent
measurement noise and correlated raw-data score noise can react differently
to a small multiplication singular value. A complex determinantal component
is not automatically a physically attainable real statistical submodel.

\subsection{A precise next target}
Let a parameter target be $\tau(\lambda)$ and let an actual feature map
be $\mathsf O(\lambda)$. Its inverse modulus is
\[
 \omega(\delta)=\sup\{d(\tau(\lambda),\tau(\lambda')):
   d_Z(\mathsf O(\lambda),\mathsf O(\lambda'))\le\delta\}.
\]
If a finite-memory procedure estimates that feature with a proved common
error bound $\delta$, fitting the model on a compact parameter set gives
an error at most $\omega(2\delta)$: the true parameter is an admissible
fit, and the triangle inequality bounds the discrepancy of the two model
features. A matched statistical lower bound requires actual distributional
closeness. If
$\norm{P_{n,\lambda_0}-P_{n,\lambda_1}}_{\TV}\le a<1$, any estimator
has maximum squared target risk at least
\[
 \frac{1-a}{8}
 d(\tau(\lambda_0),\tau(\lambda_1))^2.
\]
Indeed choose the nearer of the two target points as a test; on an incorrect
choice the estimation error is at least half their separation. The two
error probabilities sum to at least $1-a$.

These two elementary bounds specify a route to a resource--statistical
singularity theorem, but do not prove a universal matched modulus law.
Covering, testing, attainable alternatives, and tail control must be supplied.
The geometry-of-rates literature, including \cite{DonohoLiu1991}, is the
appropriate comparison, not a claim that all minimax theory follows from
one rank locus.

\subsection{The physical realization obligations remain}
For deterministic hyperbolic and particle systems, one must still establish
the true experiment kernels, physical clock handling, tightness and
recovery at exponential scale, and the domains of any response or memory
operators. A pathwise finite-state prediction error does not provide a
Fourier local limit theorem. A static covariance factorization does not
construct a collision process. A projected state does not acquire a Markov
semigroup merely because it is named a predictive coordinate.

The historical typed-contraction program \cite{QiC2} can therefore be
reorganized as a family of verified realization theorems attached to these
interfaces. Its model-specific proofs are not assumed completed by the
present definitions. In particular memory response in a varying Hilbert
or form domain requires its own transport and domain arguments.

\subsection{Conclusion}
A common foundation is obtained by keeping the entire experiment ahead of
its geometric representations. Positive preparation and future tests give
prediction; a real acquisition law gives a distortion problem; stable
updates give an online implementation; a causal simulator gives a
comparison, with its state and error charged. In the refreshing model these
steps yield a uniform finite-memory law through loss of the acquisition
channel. The separate entropy, score, and exponential identities explain
which forms of information transport are valid. Stronger singular and
long-time results must enter through corresponding theorems, not through
additional meanings assigned to the word ``contraction''.
